\section{Introduction}
\label{sec:intro}

Epistemic uncertainty is said to be what more data would reduce, but the phrase leaves open
\emph{which} data. Most deployments can buy three kinds: more unlabelled inputs from the
deployment stream, more labelled examples of the classes already known, and labelled examples of
a class not seen before. A single scalar does not say which of these would help at a given
input. We make the acquisition explicit for a setting practitioners already
have: a classifier trained once and then frozen, the labelled data it was trained against, and a
growing unlabelled log of the inputs it meets in deployment. From the log we form the set of
deployment distributions that are mixtures of the known classes (plus one candidate novel source)
and lie within a computable distance of what the log shows, and the spread of the label
posteriors this set admits splits into three parts: prevalence, label, and coverage of novelty.

The split matters because the parts answer to different acquisitions. When the class distributions
are known, uncertainty about how common each class is falls only with more of the deployment stream,
uncertainty about the label rates falls only with labels, and a novel source can be told apart from a
change in known prevalences only when its inputs fall somewhere the known classes do not. When the
class distributions are themselves estimated from labels, as in practice, labels and log both move the
prevalence width, at an exchange rate the theory motivates: each acquisition acts through the part of
the radius it controls. A score that merges everything, such as ensemble disagreement, moves
with every acquisition that feeds it and cannot say which one helps (\cref{sec:exp-synthetic}).
\citet{young2026epistemic} argues the conceptual point, that reducibility is relative to an
acquisition class; we supply estimators, finite-sample inclusion and tests of the acquisition
decisions they imply, including one that fails.

\paragraph{Approach.} The frozen classifier maps inputs to finitely many cells, so each class is a
distribution over cells and the deployment law is a mixture of them. With the categorical kernel,
the maximum mean discrepancy (MMD) between such laws is the Euclidean distance between cell
histograms, so the set of mixtures consistent with the log is a ball whose radius has an explicit
constant. Each width is the range of one posterior over that ball, computed without relaxation by a
second-order cone program. When the class distributions are themselves estimated from labels,
we widen the intervals to range over their estimation error as well, and these joint intervals
carry the coverage claim.

\paragraph{Contributions.}
\begin{itemize}
\item Three widths with distinct estimands (prevalence, label, coverage) and joint intervals that keep
finite-sample inclusion when class distributions are estimated, including a tight version that counts
reference error once (\cref{sec:setup}).
\item Conditions under which each width contracts, the exchange rate between log and labels they imply,
and where coverage identification stops, with machine-checked cores (\cref{sec:theory}).
\item A pre-registered synthetic study that varies log size, labeled sample size and label noise: with
known rows each width responds to the acquisition it tracks, ensemble mutual information responds to all
three, the ball includes the truth in every seed, and EMQ with a bootstrap loses inclusion where
prevalence is not identified (\cref{sec:exp-synthetic}).
\item A CIFAR-10 audit with two frozen encoders against a finite-sample constraint region and a
two-stage bootstrap (\cref{sec:exp-real}); a registered test showing that under shift the widths do not
choose data better than a constant policy (\cref{sec:exp-blindspot}); and a registered test on a blind
encoder showing that, when the class-conditionals hold, the radius exchange rate picks the acquisition
that tightens a certified interval more; on real features the gain over always buying more log is small,
the choice is per deployment rather than per input, and data-free bookkeeping of the two radii does as
well as the full method (\cref{sec:exp-exchange}).
\item A per-input attribution scored against an oracle: the argmax of the widths fails, and a registered
pinned-estimate rule names the limiting source better than any input-independent rule on CIFAR-10,
Gaussian-feature and synthetic worlds; but neither it nor a registered class-wise label allocation built on
it improves an acquisition outcome over simple constants (\cref{sec:exp-attribution}). All registered
outcomes, failures included, are in \cref{app:outcomes}.
\end{itemize}
